\documentclass[11pt,twocolumn]{article}
\usepackage[a4paper,top=15mm,bottom=17mm,left=15mm,right=15mm,columnsep=7mm]{geometry}
\usepackage{kotex}
\usepackage{fontspec}
\setmainfont{NanumGothic}
\setsansfont{NanumGothic}
\usepackage{amsmath,amssymb,booktabs,enumitem,microtype,xcolor,hyperref,balance}
\hypersetup{colorlinks=true,linkcolor=black,citecolor=black,urlcolor=blue!55!black}
\setlength{\parindent}{0.8em}
\setlength{\parskip}{0pt}
\setlist{nosep,leftmargin=1.4em}
\renewcommand{\baselinestretch}{0.98}
\setlength{\emergencystretch}{1.2em}
\newcommand{\draftnote}[1]{\textcolor{red!65!black}{[초안: #1]}}
\renewcommand{\abstractname}{초록}
\renewcommand{\refname}{참고문헌}

\title{\vspace{-8mm}\textbf{인과 구조와 역강화학습을 결합한\\
이커머스 사용자 행동 시뮬레이터 설계}\vspace{-2mm}}
\author{익명 저자 \quad \textit{KSC 투고용 연구계획 초안}}
\date{}

\begin{document}
\maketitle
\vspace{-6mm}

\begin{abstract}
본 연구는 AmExpert 2019의 고객 거래 및 쿠폰 캠페인 이력을 다변량 시계열로 재구성하고,
NTS-NOTEARS로 비선형 동적 인과 그래프를 학습한 뒤, 그래프 기반 상태 특징과 최대엔트로피
역강화학습(MaxEnt IRL)을 결합하여 사용자 효용을 복원하는 합성 이커머스 환경을 제안한다.
사용자별 효용계수는 혼합 가우시안으로 계층화하고, 온도 보정 다항로짓 선택모형으로
부분적으로 결정적이면서 확률적인 행동을 생성한다. 생성 환경에서는 Bandit, PPO, SAC와
GRU4Rec을 동일한 관측·행동 공간에서 비교한다. 본 초안은 데이터 전처리, 식별 가정,
모형 구성 및 검증 계획을 명시하며, 실험 결과는 후속 구현 후 채울 예정이다.
\end{abstract}
\noindent\textbf{주제어:} 인과발견, NTS-NOTEARS, 역강화학습, 선택모형, 사용자 시뮬레이션

\section{서론}
추천·프로모션 정책의 오프라인 평가는 로그에 존재하지 않는 행동의 반응을 관찰할 수 없다는
근본적 한계를 갖는다. 단순 재표본화는 사용자 상태 변화와 정책 개입의 방향성을 보존하지
못하고, 예측모형만으로 만든 시뮬레이터는 상관관계를 개입 효과로 오해할 수 있다. 이에 본
연구는 (1) 시간 지연과 동시점 관계를 포함하는 동적 인과 구조, (2) 관찰 행동으로부터 복원한
사용자 효용, (3) 고객 이질성과 확률적 선택을 하나의 생성과정으로 결합한다.

핵심 연구질문은 다음과 같다. 첫째, 캠페인과 구매가 중첩된 패널에서 안정적인 Markov
boundary를 발견할 수 있는가? 둘째, 이 특징 집합으로 추정한 보상은 관찰 행동과 구매 반응을
재현하는가? 셋째, 합성 환경에서 우수한 정책 순위가 관측자료의 보유기간 평가에서도 유지되는가?
본 연구의 기여는 단일 DAG를 정답으로 선언하는 데 있지 않다. 부트스트랩과 시간해상도 변화에도
남는 구조만 생성기의 제약으로 사용하고, 알려진 DAG를 가진 합성 실험을 먼저 통과시키는
검증 중심 파이프라인을 제안한다.

\section{데이터와 전처리}
\subsection{AmExpert 2019의 재구성}
AmExpert 2019는 고객 인구통계, 상품 메타데이터, 일자별 구매, 쿠폰--상품 매핑, 캠페인 기간,
고객별 쿠폰 수신과 사용 여부를 제공한다. Langen과 Huber~\cite{langen2023}에 따르면 원자료는
2012년 1월부터 2013년 7월까지의 구매와 시간적으로 일부 중첩된 18개 캠페인을 포함한다.
원 논문은 쿠폰 유효 집합이 기간 안에서 일정하도록 모든 시작·종료 경계에서 시간을 다시 나누어
33개 인공 캠페인 기간을 만들었다. 기간별 결과는 고객의 총지출을 유효일수로 나눈 일평균 지출이며,
각 고객을 33개 기간에 걸쳐 결합하였다. 통제변수는 인구통계, 기간 고정효과, 상품군별 사전
일평균 지출, 직전 기간의 쿠폰 수신·사용, 그리고 관심 상품군 외 동시 노출 쿠폰 더미이다.

원 논문의 전체 표본은 50,624개 고객--기간 관측치이며, 15,327개(약 30\%)가 하나 이상의
쿠폰을 수신했다. 수신 관측치 중 실제 사용률은 3\%에 불과하다. 연령·가구규모·주거형태·소득의
\texttt{unknown}은 각각 52\%, 혼인상태는 72.8\%였으며, 논문은 이를 별도 범주로 유지한 주 분석과
인구통계가 모두 관측된 431명(13,792개 관측치)의 완전사례 강건성 분석을 병행하였다. 따라서
쿠폰사용 희소성뿐 아니라 정보 관측 여부와 타기팅의 연관성도 검증해야 한다.

위 절차는 원 논문의 효과추정용 전처리이다. 이와 구분하여, \textbf{본 연구가 새로 제안하는}
NTS-NOTEARS 입력의 주 분석 단위는 고정 간격인 \textbf{고객--주(customer-week)}이다.
고객 $i$, 주 $t$에 대해
\[
\begin{aligned}
x_{it}=(&{\rm spend,qty,items,categories},\\[-1mm]
        &{\rm discount,offer,redeem})
\end{aligned}
\]
를 구성한다. 지출·수량·고유 상품/상품군 수, 일반·쿠폰 할인액, 신규 및 활성 쿠폰 수,
보고된 사용 여부를 포함한다. 인구통계의 결측은 삭제하지 않고 ``unknown'' 범주와 관측 마스크로
분리한다. 원자료에 정확한 쿠폰 발송·사용 시각이 없으므로 발송은 캠페인 시작 주, 사용 보고는
종료 주에 배치한다. 이는 원 논문의 처리법이 아니라 본 연구의 측정 가정이며, 시작·종료 주
배치를 바꾼 민감도 분석을 함께 수행한다. 거래가 없는 주도 0으로 채워 균일한 간격을 유지하되,
구조적 0과 결측을 구분한다.

고객별 시계열은 독립된 배열로 보존한다. 마지막 고객의 마지막 주와 다음 고객의 첫 주를
연결하면 거짓 lag edge가 생기므로, lag window는 고객 경계를 절대 넘지 않는다. 고객 단위로
학습/검증/시험을 분리하고, 시간 누수를 막기 위해 각 고객 내부에서도 과거에서 미래 순으로
절단한다. 원 논문의 33개 캠페인 기간 패널은 동일한 일평균 정의로 별도 재현하여 주 단위 결과의
시간 집계 민감도 분석에 사용한다.

\section{제안 방법}
\subsection{NTS-NOTEARS 기반 동적 인과발견}
NTS-NOTEARS~\cite{sun2023}는 1차원 CNN으로 비선형 조건부 함수를 근사하며 lagged
(inter-slice) 관계와 instantaneous(intra-slice) 관계를 함께 학습한다. 시차 $L\in\{2,4\}$에
대해 각 변수 $X_j^t$를 $X^{t-L:t}$로 예측하고 다음 목적함수를 최적화한다.
\begin{align*}
\min_{\theta,W}\quad &\frac{1}{n}\sum_t\|X^t-f_\theta(X^{t-L:t})\|_2^2
+\lambda_1\|W\|_1,\\[-1mm]
\text{s.t.}\quad &h(W^{(0)})=0,
\end{align*}
\[
h(W^{(0)})=\operatorname{tr}\!\left[\exp(W^{(0)}\circ W^{(0)})\right]-d.
\]
동시점 그래프 $W^{(0)}$에만 비순환 제약을 부여하고 과거에서 미래로의 edge만 허용한다.
캠페인 시작 시 고정된 \texttt{offers\_started}로 향하는 동시점 구매 edge는 금지하되, 과거
구매에서 미래 offer로의 edge는 관측 타기팅을 반영하므로 허용한다. 고객 부트스트랩, lag,
임계값 변화에서 선택확률이 $\pi_0$ 이상인 edge만 안정 그래프에 남긴다. DYNOTEARS와
PCMCI+를 구조 민감도 기준선으로 둔다~\cite{pamfil2020,runge2019}.

\subsection{그래프 기반 특징 추출과 MaxEnt IRL}
구매·지출·쿠폰사용 노드 $Y$에 대해 안정 그래프의 부모, 자식 및 자식의 부모로 구성되는
Markov boundary 후보 $\mathcal{M}(Y)$를 구한다. 단일 그래프의 오류 전파를 줄이기 위해
edge 선택확률, lag, 방향을 특징에 함께 부여하고, 정책결정 이후 관측되는 변수는 상태에서
제외한다. 최종 상태 $s_t$는 안정 특징, 최근 $L$주 요약, 정적 고객정보 및 노출 가능한
행동 집합으로 구성한다.

쿠폰 미수신을 포함한 캠페인 행동 $a_t$가 선형 보상
$r_{\beta_u}(s_t,a_t)=\beta_u^\top\phi(s_t,a_t)$를 최대화한다고 가정한다. MaxEnt IRL은
\begin{align*}
p(\tau_u\mid\beta_u)&\propto
\exp\!\left\{\sum_t\beta_u^\top\phi(s_t,a_t)\right\},\\
\hat\beta&=\arg\max_\beta\sum_u\log p(\tau_u\mid\beta)-\lambda\|\beta\|_2^2
\end{align*}
를 이용해 관측 궤적과 모형의 기대 feature count를 맞춘다~\cite{ziebart2008}. 고객별 로그가
짧으므로 완전 독립 추정보다 전역 평균에 shrinkage되는 계층형 $\beta_u$를 사용한다. 정책의
support가 부족한 행동에 대한 보상은 식별되지 않으므로 propensity와 유효표본크기(ESS)를
함께 보고한다. 단, AmExpert에는 진열·재고를 포함한 완전한 소비자 선택집합이 없으므로
$\beta_u$는 상품 수준의 절대 선호가 아니라 관측된 쿠폰 범주와 구매 요약에 대한 상대 효용으로
해석한다.

\subsection{GMM 이질성과 다항로짓 선택}
추정된 계수와 정적 문맥 $z_u=[\beta_u,c_u]$에 혼합 가우시안을 적합한다.
\[
p(z_u)=\sum_{k=1}^{K}\pi_k\mathcal{N}(z_u\mid\mu_k,\Sigma_k),
\]
여기서 $K$는 BIC와 부트스트랩 군집 안정도로 정한다. 새 사용자는 군집을 먼저 뽑고
$\beta_u$를 조건부 표본추출하여 군집 내·군집 간 이질성을 모두 보존한다.

행동은 외부선택(outside option)을 포함한 다항로짓으로 생성한다~\cite{mcfadden1974}.
\begin{align*}
P(a\mid s,u)&=\frac{\exp\{V_{u}(s,a)/T\}}
{\sum_{a'\in\mathcal{A}(s)}\exp\{V_u(s,a')/T\}},\\[-1mm]
V_u(s,a)&=\beta_u^\top\phi(s,a).
\end{align*}
온도 $T$가 작을수록 결정적이고 클수록 무작위에 가깝다. $\beta/T$는 별도로 식별되지 않으므로
$\|\beta\|$ 또는 $T$ 중 하나를 고정하고, 검증자료의 선택 엔트로피, top-1 비율과 상품군 점유율을
동시에 맞춘다. 확률을 임의로 sharpen하지 않고, $T$의 허용구간과 outside-option 상수를
사전등록해 과도한 결정성과 과도한 무작위성을 모두 방지한다.

\section{실험 및 평가 계획}
첫째, 알려진 동적 DAG와 잠재 사용자 군집을 가진 반합성자료에서 edge F1, SHD, 방향 및 lag
정확도를 측정한다. 둘째, AmExpert에서 고객 부트스트랩 edge 안정도, 보유기간 예측오차,
행동확률 calibration과 생성자료의 marginal/transition distance를 평가한다. 셋째, 생성기에
개입하여 쿠폰 강도·상품군 정책의 dose-response와 방향 일관성을 점검하되, 관측 AmExpert만으로
실세계 개입효과가 식별되었다고 주장하지 않는다.

정책 비교군은 무작위·인기·contextual bandit, PPO~\cite{schulman2017}, SAC~\cite{haarnoja2018},
GRU4Rec~\cite{hidasi2016}이다. SAC는 연속 할인강도, PPO는 이산/혼합 쿠폰 정책에 적용하고,
GRU4Rec은 다음 상품군 또는 쿠폰 순위 예측기로 사용한다. 공통 후보집합과 예산 아래 누적 효용,
구매·사용률, 매출, 다양성, 장기 유지, regret 및 정책 간 순위 안정성을 보고한다. simulator가
특정 학습기에게 유리하지 않은지 확인하기 위해 구조, 보상, choice temperature를 각각
교란하는 ablation을 수행한다.

\section{예상 기여와 한계}
제안 파이프라인은 ``상관 기반 로그 $\rightarrow$ 인과 가설 $\rightarrow$ 효용 분포
$\rightarrow$ 확률적 선택 $\rightarrow$ 정책평가''를 재현 가능한 단위로 연결한다. 그러나
AmExpert의 쿠폰 할당은 무작위화가 확인되지 않았고, 정확한 발송·사용 시점, 비구매 시 고려한
상품, 진열·재고와 경쟁상품 노출이 없다. 선택-관측가능성 가정이 위배되면 구조와 보상이 함께
편향된다. 따라서 학습 DAG는 인과 정답이 아니라 검증 대상 가설 그래프이며, MaxEnt IRL 보상도
관측정책과 구성한 선택집합 아래의 상대 효용이다. 최종 논문에서는 외부 randomized uplift
자료를 보정·검증에 추가하거나, 불가능할 경우 부분식별 범위와 이 한계를 결과 해석의 전제로
둔다. \draftnote{실험 수치, 표본 생성 규칙과 최종 하이퍼파라미터는 구현 후 삽입한다.}

\balance
\begingroup
\footnotesize
\setlength{\parindent}{0pt}
\begin{thebibliography}{10}
\bibitem{langen2023} H. Langen and M. Huber, ``How causal machine learning can leverage marketing strategies: Assessing and improving the performance of a coupon campaign,'' \textit{PLoS ONE}, 18(1):e0278937, 2023.
\bibitem{sun2023} X. Sun, O. Schulte, G. Liu, and P. Poupart, ``NTS-NOTEARS: Learning Nonparametric DBNs With Prior Knowledge,'' \textit{AISTATS}, pp. 1942--1964, 2023.
\bibitem{zheng2018} X. Zheng et al., ``DAGs with NO TEARS: Continuous Optimization for Structure Learning,'' \textit{NeurIPS}, 2018.
\bibitem{pamfil2020} R. Pamfil et al., ``DYNOTEARS: Structure Learning from Time-Series Data,'' \textit{AISTATS}, 2020.
\bibitem{runge2019} J. Runge et al., ``Detecting and Quantifying Causal Associations in Large Nonlinear Time Series Datasets,'' \textit{Science Advances}, 5(11), 2019.
\bibitem{ziebart2008} B. Ziebart et al., ``Maximum Entropy Inverse Reinforcement Learning,'' \textit{AAAI}, 2008.
\bibitem{mcfadden1974} D. McFadden, ``Conditional Logit Analysis of Qualitative Choice Behavior,'' in \textit{Frontiers in Econometrics}, 1974.
\bibitem{schulman2017} J. Schulman et al., ``Proximal Policy Optimization Algorithms,'' arXiv:1707.06347, 2017.
\bibitem{haarnoja2018} T. Haarnoja et al., ``Soft Actor-Critic: Off-Policy Maximum Entropy Deep Reinforcement Learning with a Stochastic Actor,'' \textit{ICML}, 2018.
\bibitem{hidasi2016} B. Hidasi et al., ``Session-based Recommendations with Recurrent Neural Networks,'' \textit{ICLR}, 2016.
\end{thebibliography}

\section*{부록 A. 재현성 및 전처리 체크리스트}
\begin{enumerate}
\item 원본 6개 테이블의 행 수, schema, checksum을 manifest로 기록하고 원자료는 배포하지 않는다.
\item 금액 부호, 중복 거래, 캠페인 날짜 형식, 상품 매핑 cardinality와 결측률을 검증한다.
\item 원 논문의 33개 기간, 일평균 지출, 사전 상품군 지출, 직전·동시 쿠폰 통제를 재현한다.
\item 고객 경계 보존 offset을 저장하고 모든 split·bootstrap을 고객 단위로 수행한다.
\item 주 단위 패널과 33개 인공 캠페인 패널, lag 2/4의 안정 edge를 비교한다.
\item unknown 범주 주 분석과 완전사례 분석을 병행하고 사용률 3\%의 불균형을 보고한다.
\item seed, forbidden edge, threshold, GMM $K$, 온도 $T$와 모든 평가 지표를 공개한다.
\end{enumerate}
\endgroup

\end{document}
